Continuous tuning of the common Bloch radius of an elegant joint measurement requires its local regular-tetrahedral geometry to remain compatible with a complete orthonormal measurement basis. This paper presents an explicit continuous family of elegant joint measurements for every integer $n\ge5$, using a unified phase construction for odd and even $n$. The family passes through a specified reference measurement, preserves orthonormality and fixed local tetrahedral directions, and continuously covers the radius interval $[0,2r_E(n)]$, where $r_E(n)=\sqrt3/2^{n-1}$ is the reference radius. Within a fixed phase representation, a complete classification is obtained of the radius laws permitted by straight phase paths with two-valued real generators that pass through the reference measurement and preserve regular tetrahedral geometry with a common radius in its neighborhood. Consequently, the bound $2r_E(n)$ is attained and optimal within this model for every $n\ge5$, and $n=5$ is the smallest system size allowing a radius increase in this model. For every even $n\ge8$, a two-dimensional measurement family is further constructed whose two coordinate sections recover the new radius-increasing family and the known cosine family column by column. The low-dimensional results further reveal the role of the phase path's shape: at $n=4$ in the same phase representation, even with arbitrary real phase generators, no straight phase path through the reference measurement that preserves the stated geometric conditions in its neighborhood can increase the radius, whereas an explicit curved phase path reaches $(1+1/\sqrt2)r_E(4)$.
For every $n\ge5$, the entire path remains genuinely multipartite entangled, with its concurrence determined exactly by the common radius and its weakest bipartitions isolating a single qubit.
